We replicate 7 tasks from \citet{gengBenchmarkingPtOPnO2024}'s end-to-end prediction and optimization (PnO) benchmark, using their implementation and settings. They use synthetic and real datasets from energy, citation networks, and beyond for a variety of decision tasks. See App.~\ref{app:bench} for task specific descriptions.

We add a DPO method to the benchmark based on \citet{TueroPytorchDPO}'s code. While \citeauthor{gengBenchmarkingPtOPnO2024}'s provided software contains some DPO code, it is not fully implemented. Results for the 11 baseline methods are taken from the original benchmark.

\textbf{Training DPO.} Our basic or ``naive'' DPO completes a grid search of noise-level $\sigma \in [0.1, 0.5, 1]$ and learning rate $[0.01, 0.005, 0.001]$, selecting the lowest-cost run on the benchmark's validation set. These 9 runs for DPO may advantage it compared to other methods (which received 5 different learning rates and no tuning of other hyperparameters). Yet our goal is to highlight what is possible with DPO, not claim superiority. 

\textbf{Results.} Figure \ref{fig:bench_bump} shows DPO's performance across the 7 tasks and 11 methods studied by \citeauthor{gengBenchmarkingPtOPnO2024}. DPO achieves the best performance on bipartite matching and budget allocation, near-optimal performance on the real-world knapsack energy problem, and ranks in the middle for 2 more tasks.
%the synthetic knapsack and energy scheduling task, and while performance on TopK is catastrophic compared to the best method, it actually outperforms several other methods.

Here, some results contrast with the findings of \citet{gengBenchmarkingPtOPnO2024}. While the authors say "methods in the discrete category [such as DPO] do not perform as well as
others, especially on non-linear optimization tasks (budget allocation)," we find that DPO with basic grid search performs the best at the Budget Alloc.~task.  
% budget allocation is truly non-linear (submodular), as is portfolio (quadratic in decision for risk minimization) but topk is absolutely linear

\textbf{Extra tuning.}
If, after the initial size-9 grid search, DPO was in the bottom half of the methods compared, we performed further tuning to illustrate what can help and harm DPO's performance. 

On the TopK task where DPO-basic scored quite poorly, our extra tuning
involved decaying $\sigma$ linearly to 0 over iterations, recommended by \citet{cordonnierDifferentiablePatchSelection2021}. We think this decay is helpful due to the nature of the problem: a decision boundary which rarely appears in the training data matters a lot in overall test cost.
For this task, the model $\hat{y}_{\theta}$ is well-specified, capable of perfectly recovering the data generating process. 
Thus, a two-stage rather than end-to-end method performs best overall, which is explained by  proofs in \citet{elmachtoubEstimateThenOptimizeIntegratedEstimationOptimizationSample2025}.


%This al
%A two-stage approach works best for this task, as the model used is well-specified and can perfectly represent the data generating process. This aligns with the proofs provided in \citet{elmachtoubEstimateThenOptimizeIntegratedEstimationOptimizationSample2025}.

The energy task exposes one of the major weaknesses of DPO: each noisy sample requires its own solver call, which is expensive here. As such, we were only able to use $M{=}1$ sample when training on this task, and we believe this to be the reason for its failure. Further tuning (more $\sigma$ values, $\sigma$ decay) were unsuccessful: we obtained better regret on the validation set, but slightly worse test-set performance.

Similarly, it was challenging to improve performance on the synthetic Knapsack task. Because $\sigma$ is either constant or decaying, it is important that the predicted outcome $y$ stays in a relatively stable range of magnitudes. As such, we placed a $\text{tanh}$ activation on the outputs of the model for this task. Again, we see this improves validation but harms test-set performance.

Finally, the portfolio task is a continuous quadratic problem, and therefore already has informative gradients and requires no alternative optimization objective. Accordingly, DPO-tuned in this problem corresponds to optimizing portfolio returns directly.